\documentclass[10pt, a4paper]{article}
\usepackage[T1]{fontenc}
\usepackage[utf8]{inputenc}
\usepackage{lmodern}
\usepackage[margin=1.8cm, top=2.3cm, bottom=2.3cm, headheight=14pt]{geometry}
\usepackage{amsmath}
\usepackage{amssymb}
\usepackage[table]{xcolor}
\usepackage{tikz}
\usetikzlibrary{calc}
\usepackage{tcolorbox}
\tcbuselibrary{breakable}
\usepackage{tabularx}
\usepackage{booktabs}
\usepackage{enumitem}
\usepackage{parskip}
\usepackage{titlesec}
\usepackage{fancyhdr}
\usepackage{hyperref}
\hypersetup{colorlinks=true, linkcolor=accent, urlcolor=accent, citecolor=accent}

\definecolor{accent}{RGB}{22,163,74}
\definecolor{accentlight}{RGB}{240,253,244}
\definecolor{accentmid}{RGB}{34,197,94}
\definecolor{errred}{RGB}{220,38,38}
\definecolor{errbg}{RGB}{254,242,242}
\definecolor{orng}{RGB}{249,115,22}
\definecolor{orngbg}{RGB}{255,247,237}
\definecolor{retreival}{RGB}{29,78,216}
\definecolor{retrievalbg}{RGB}{239,246,255}
\definecolor{slate}{RGB}{71,85,105}
\definecolor{slatelight}{RGB}{148,163,184}
\definecolor{dblue}{RGB}{29,78,216}

\newcommand{\dgr}{\ensuremath{^\circ}}

\tcbset{
  note/.style={colback=accentlight,colframe=accentmid,
    left=6pt,right=6pt,top=4pt,bottom=4pt,breakable},
  artnote/.style={colback=orngbg,colframe=orng,
    left=6pt,right=6pt,top=4pt,bottom=4pt,breakable}
}

\titleformat{\section}{\large\bfseries\color{accent}}{\thesection}{0.6em}{}[\color{accentmid}\titlerule]
\titleformat{\subsection}{\normalsize\bfseries\color{slate}}{}{0em}{\MakeUppercase}
\titlespacing{\section}{0pt}{16pt}{8pt}
\titlespacing{\subsection}{0pt}{12pt}{4pt}

\pagestyle{fancy}
\fancyhf{}
\fancyhead[L]{\small\color{slate}STEAM, Grades 5--9 --- Course Context}
\fancyhead[R]{\small\color{slate}Reference Sheet}
\fancyfoot[C]{\small\color{slatelight}\thepage}
\renewcommand{\headrulewidth}{0.4pt}
\renewcommand{\headrule}{\color{slatelight}\hrule width\headwidth height\headrulewidth}

\setlist[itemize]{leftmargin=1.4em, itemsep=3pt, topsep=2pt}
\setlist[enumerate]{leftmargin=1.6em, itemsep=4pt, topsep=2pt}

\begin{document}

\begin{tcolorbox}[colback=accentlight, colframe=accent, leftrule=5pt,
    toprule=0.6pt, bottomrule=0.6pt, rightrule=0.6pt,
    top=8pt, bottom=8pt, left=12pt]
  {\small\bfseries\color{accent}COURSE CONTEXT \textperiodcentered{} SUPPLEMENTARY PROGRAM}\\[4pt]
  {\LARGE\bfseries STEAM, Grades 5--9}\\[3pt]
  {\small\color{slate}Read alongside \texttt{\_meta-context.md} and \texttt{research/steam-progression-research.md}.}
\end{tcolorbox}
\vspace{10pt}

\begin{tabularx}{\linewidth}{@{}l X l X@{}}
\textbf{Grades} & 5--9, one year plan per grade & \textbf{Year length} & 24 weeks \\
\textbf{Schedule} & 2-hour session every 3 days & \textbf{Hours / grade} & $\sim$40 sessions, $\sim$80 hrs \\
\end{tabularx}
\vspace{4pt}

\section{Course Overview}

\textbf{Session count assumption} (state this before using the plans): ``Every 3 days'' is read as roughly 1.6--1.7 sessions/week on a 5-day school week --- \textbf{$\sim$40 sessions/year, 2 hrs each, $\sim$80 instructional hours per grade.} If the real rotation (e.g.\ a lettered-day timetable) produces a different count, unit weightings scale proportionally --- the ratio between units matters more than the absolute session count. Flag it if the real number falls under 32 or over 48 sessions; the per-grade plans should be re-paced.

\textbf{Audience:} students and teachers. This is a supplementary program, not a mandated Alberta strand --- see the Alberta angle in the research file. Materials should still be print/deliver-ready for any teacher picking up the sequence.

\section{Evidence Foundation}

Every structural decision in this course traces to \texttt{research/steam-progression-research.md}. Read it for the full ``why.'' The short version, restated as course rules:

\begin{enumerate}
\item \textbf{Concrete before abstract}, faded one step per grade band --- not skipped, not held onto too long.
\item \textbf{All five strands run every unit, every year} --- coding, robotics, spatial/construction, art/design, and the engineering design cycle are never taught as separate terms.
\item \textbf{The design cycle recurs with escalating constraints}, not as a one-off unit.
\item \textbf{At least one open-ended task per unit}, even in Grade 5 --- guided-build only produces instruction-followers.
\item \textbf{Every unit ends with a share-and-reflect step} --- the artifact is not the finish line, explaining it is. Capstone units (Unit 4) end with a produced video, not just a live demo.
\end{enumerate}

\begin{tcbraster}[raster columns=1, raster equal height=false]
\begin{tcolorbox}[artnote, title={\small\bfseries\color{orng}On the name ``STEAM''}]
This is not STEM with robots bolted on. The A is a required strand with its own dedicated tasks, not a paint job applied to robotics builds --- see \S3 Art \& Media Strand. Every unit needs at least one task that does not involve a robot at all.
\end{tcolorbox}
\end{tcbraster}

\section{Equipment Map}

The platforms on hand map onto the concrete-to-abstract fade almost exactly as the research recommends --- the backbone of the multi-year sequence. Each grade's ``primary'' platform is new; prior platforms return as the \emph{sensor/logic} layer under a more complex build, not discarded.

\begin{center}
\small
\begin{tabularx}{\linewidth}{@{}p{2.7cm}p{2.4cm}Xp{2.1cm}@{}}
\toprule
\textbf{Platform} & \textbf{Coding mode} & \textbf{Strongest use} & \textbf{Primary grade} \\
\midrule
Dash (Wonder Workshop) & Blockly app, very tangible & First exposure to sequence/loop/conditional; low floor for reluctant starters & Grade 5 (retired after; concepts carry forward) \\
micro:bit & MakeCode blocks $\rightarrow$ Python & Physical computing, sensors (accelerometer, light, radio), wearables & Grade 5--6 intro, Gr.\,7+ as sensor module \\
Sphero & Sphero Edu Blockly $\rightarrow$ JavaScript & Path/maze programming, first block-to-text bridge & Grade 6 intro, Gr.\,7 text transition \\
LEGO Spike Prime & Word Blocks $\rightarrow$ Python & Build + code integration; motors, multiple sensors & Grade 7 intro, Gr.\,8--9 primary \\
REV kits (IQ-class) & Paired with Spike/micro:bit logic & Competition-grade construction, structural rigor & Grade 8 intro, Gr.\,9 capstone \\
Maker/construction materials & N/A --- physical only & Spatial reasoning, mechanism literacy before motorizing & Every grade, every unit \\
3D printers (MakerBot Sketch) + Tinkercad & Not coding --- digital design/fab & Turns Build Challenge/Capstone into a real design-print-test-redesign loop & Gr.\,6 intro, Gr.\,7 CAD gap, Gr.\,8--9 custom parts \\
\bottomrule
\end{tabularx}
\end{center}

\noindent\textbf{Reading the table as a coding-abstraction ladder:} Dash (fully tangible) $\rightarrow$ micro:bit blocks $\rightarrow$ Sphero blocks $\rightarrow$ Sphero JavaScript / Spike Prime blocks $\rightarrow$ Spike Prime Python. No grade should jump more than one rung in a single unit.

\section{Art \& Media Strand}

The equipment map is a coding/robotics ladder because that's what needs platform sequencing across five years. Art doesn't need new hardware to fade the same way --- it needs a standing requirement that keeps it from being displaced by whichever robot is newest. Two rules do that.

\subsection{Rule 1 --- a dedicated non-robot art/design task, every unit}

Not aesthetics applied to a robot build (a nicer paint job on the Dash attachment doesn't count) --- a task whose deliverable is a drawing, physical model, diagram, storyboard, or other visual/design artifact that stands on its own, produced with no robot involved. This can overlap with the ``at least one open-ended task per unit'' rule but doesn't have to.

\begin{itemize}
\item \textbf{Unit 1 (Foundations):} a sketch/diagram task --- students draw and label the program flow they're about to build, or sketch the object/character a robot behavior represents, before touching a device.
\item \textbf{Unit 2 (Sensors \& Response):} a visual model of the environment the robot responds to --- a labeled map, diagram, or scale drawing of the test environment, made before the live test.
\item \textbf{Unit 3 (Build Challenge):} the design-review artifact (sketch + material list) counts, treated as a real drawn/rendered plan, not a rough note --- drafting skill is most directly on display here.
\item \textbf{Unit 4 (Capstone):} the video itself is the primary art/media deliverable --- storyboarding and shooting/narrating it is design work, not paperwork tacked onto the end.
\end{itemize}

\subsection{Rule 2 --- Unit 4 ends with a produced video}

Every grade's Unit 4 (Capstone) ends with a produced video, not a live-only demo. Storyboard, planning, filming or screen-capturing the build/program in action, and narrating the design decisions and iteration history --- replacing, or where a live audience is available, supplementing the verbal-only share-and-reflect used in Units 1--3.

This is deliberate scope: Units 1--3 stay lighter-touch so the video doesn't become routine busywork; it's reserved for the moment in the cycle with a real design story worth documenting. Video complexity fades the same way coding mode does --- a Grade 5 video can be a simple phone/tablet recording with spoken narration; by Grade 9, students handle their own shot planning, editing, and narration with minimal scaffolding.

\begin{tcolorbox}[note, title={\small\bfseries\color{accent}Digital fabrication note --- this changes pacing, not just materials}]
A print job takes 30 minutes to several hours. No unit can plan a design-print-test loop inside a single 2-hour session: build at least one session's buffer between ``design finalized'' and ``part in hand,'' run print queues overnight or between sessions, and give students a non-printing task while a print runs. With a small number of MakerBot Sketch units and a full class needing parts, stagger submissions across the unit rather than expecting everyone's part the same day.
\end{tcolorbox}

\section{Unit Archetype}

Same shape, every grade, every year. Each year is four units of $\sim$10 sessions ($\sim$20 hrs, $\sim$6 weeks). The unit \emph{shape} stays constant so students recognize the pattern by Grade 7 and run it with less scaffolding --- what changes year to year is the platform, the constraint complexity, and how much of the process the student initiates independently.

\begin{center}
\small
\begin{tabularx}{\linewidth}{@{}p{2.5cm}Xp{2.3cm}X@{}}
\toprule
\textbf{Unit} & \textbf{Focus} & \textbf{Constraint level} & \textbf{Ends with} \\
\midrule
1. Foundations & Re-establish/introduce this year's primary platform; name the CT vocabulary explicitly (sequence, loop, conditional, variable, debug) & Low --- single clear goal, materials given & Working artifact + share-and-reflect \\
2. Sensors \& Response & Programming against the physical world --- line-following, obstacle response, light/sound/motion triggers & Medium --- stated environment to respond to & Robot/program reacts correctly to a live class test \\
3. Build Challenge & Construction- and spatial-reasoning-heavy; mechanism literacy tied to the platform's actuators & Medium-high --- competing constraints (speed vs.\ stability, weight limit) & Design review --- peer critique against stated constraints \\
4. Capstone Design Challenge & Open-ended brief, multiple valid solutions, full design cycle with minimal prompting & High --- real external constraint (budget, material, HS-modeled rubric) & Public share-out + documented iteration history, captured as a produced video \\
\bottomrule
\end{tabularx}
\end{center}

\noindent This table is the \emph{shape}; the grade files fill in the platform, the specific brief, and the LOs.

\section{Grade-by-Grade Summary}

\begin{center}
\small
\begin{tabularx}{\linewidth}{@{}p{0.8cm}p{2.5cm}p{2.5cm}X@{}}
\toprule
\textbf{Gr.} & \textbf{Primary platform} & \textbf{Coding mode} & \textbf{Capstone shape} \\
\midrule
5 & Dash + micro:bit (intro) & Blockly (Dash), MakeCode blocks & Guided-but-open: a Dash ``helper'' for a stated classroom problem \\
6 & micro:bit + Sphero (intro) & MakeCode blocks (deeper), Sphero Edu Blockly & Sphero maze/obstacle course, student-designed, written constraint set \\
7 & Sphero (text) + Spike Prime (intro) & Sphero JavaScript, Spike Word Blocks & Open brief: build + program Spike Prime for a task, competing constraints \\
8 & Spike Prime (primary) + REV (intro) & Spike Prime Python (scaffolded) & Peer-reviewed design-build-test-redesign with an external constraint \\
9 & REV (capstone) + Spike Prime Python & Python default; blocks as fallback only & Independent capstone mirroring an HS CTS/CTF outcome; documented iteration history required \\
\bottomrule
\end{tabularx}
\end{center}

\section{Cross-Cutting Rules}

For anyone building lessons/units from this plan.

\begin{enumerate}
\item \textbf{Never let a unit become ``coding week'' then ``building week.''} Every session should touch both, even if the ratio shifts (Unit 3 leans construction-heavy, but the platform's code should still run something by unit's end).
\item \textbf{Not everything needs to be done on a robot.} Every unit needs its dedicated non-robot art/design task --- check for it explicitly before finalizing a unit; it's easy to let robotics crowd it out.
\item \textbf{Name the CT vocabulary out loud at first use each year}, even if it's a repeat. Naming beats incidental exposure at every grade, not just the youngest.
\item \textbf{At least one task per unit has more than one correct answer.} If a unit's session plan is 100\% ``follow these build steps,'' add or swap in an open task before finalizing.
\item \textbf{The share-and-reflect step is not optional and not ``if there's time.''} Build it into the session count for the last day of every unit. Unit 4 includes video production time.
\item \textbf{Escalate one constraint dimension at a time.} Don't introduce a budget limit, a new platform, and a new sensor type in the same unit --- pick the one escalation this unit is teaching.
\end{enumerate}

\section{Open Items}

\begin{itemize}
\item Rubric design for open-ended build/capstone tasks is not yet built --- flagged in the research file as a separate need. Build before the Grade 5 Unit 4 capstone is delivered.
\item \textbf{The Art \& Media Strand and Capstone video requirement are new} --- none of the 5 grade-year files reflect them yet. Each grade file's Units 1--4 need: (a) the non-robot art/design task named explicitly per unit, and (b) Unit 4 revised to specify the video deliverable (equipment, expected length/complexity, session fit). Follow-up pass, grade by grade.
\item Video equipment/software not yet confirmed (tablets on hand for coding apps may suffice for recording; editing tool unconfirmed). Confirm before writing Grade 8--9 Unit 4, where student-led editing is assumed.
\item \textbf{Confirm the CAD tool paired with MakerBot Sketch.} Plans assume Tinkercad (standard free classroom pairing) --- correct if the school has a different tool set up.
\item Safety/tool-use procedures for hand tools in the construction strand (Grade 5+) not yet written --- needed before Unit 3 of any grade is delivered to students.
\item \textbf{3D printer capacity is not yet known} (unit count, typical print time). Grade plans assume queuing is needed; confirm actual capacity so Grade 7--9 units can be paced against real parts-per-week.
\end{itemize}

\end{document}
